% !PRINTED 456
\weberpage{456}

%<<CONT
From this we obtain:

\[
\begin{aligned}
\tfrac{1}{6} F''(y_1, y_1) &= x_1 x_3 \alpha^2 + x_2 x_1 \beta^2 + x_3 x_2 \gamma^2 \\
&\quad + x_3^2 \beta\gamma + x_1^2 \alpha\gamma + x_2^2 \alpha\beta, \\
\tfrac{1}{3} F'''(y_2, y_3) &= 2x_1 x_3 \beta\gamma + 2x_2 x_1 \alpha\gamma + 2x_3 x_2 \alpha\beta \\
&\quad + x_3^2 (\beta\gamma + \alpha^2) + x_1^2 (\alpha\gamma + \beta^2) + x_2^2 (\alpha\beta + \gamma^2).
\end{aligned}
\]
% [sic] mixed term printed with triple prime F''' here against F'' just below for the same pairing

Now since the solutions of the system (3) are of the same form ($y_1 = \alpha x_1 + \beta x_2 + \gamma x_3, \dots$), one obtains from this:

\[
\begin{aligned}
\tfrac{1}{6} F''(y_1, y_1) - y_1 y_3 &= \\
x_1 x_3 (\alpha^2 - \alpha\beta - \gamma^2) + x_2 x_1 (\beta^2 - \beta\gamma - \alpha^2) &+ x_2 x_3 (\gamma^2 - \beta^2 - \alpha\gamma), \\
\tfrac{1}{3} F''(y_2, y_3) - y_3^2 &= \\
x_1^2 (\beta^2 + \alpha\gamma - \gamma^2) + x_2^2 (\gamma^2 + \alpha\beta - \alpha^2) &+ x_3^2 (\alpha^2 + \beta\gamma - \beta^2).
\end{aligned}
\]

But the coefficients on the right-hand side of these equations $(\alpha^2 - \alpha\beta - \gamma^2) \dots$ turn out, from the values of $\alpha, \beta, \gamma$ [\S.\,115, (11)], to vanish, and so, if we further apply a cyclic permutation of the $x$, which entails a cyclic permutation of the $y$, we obtain:

\[
\begin{aligned}
&\tfrac{1}{6} F'''(y_1, y_1) = y_1 y_3, \quad \tfrac{1}{6} F'''(y_2, y_2) = y_2 y_1, \quad \tfrac{1}{6} F'''(y_3, y_3) = y_3 y_2, \\
&\tfrac{1}{3} F''(y_2, y_3) = y_3^2, \quad \tfrac{1}{3} F''(y_3, y_1) = y_1^2, \quad \tfrac{1}{3} F''(y_1, y_2) = y_2^2.
\end{aligned}
\]
% [sic] prime-count on F wavers between the two lines above, reproduced as printed

Accordingly, by \gesperrt{Euler}'s theorem [Bd.\ I, \S.\,17, (6)], we obtain:

\[
F\,(y_1, y_2, y_3) = y_1^3 y_3 + y_2^3 y_1 + y_3^3 y_2,
\]

and with this it is proved that the function $f\,(x_1, x_2, x_3)$ is in fact an invariant of our group $G_{168}$. Apart from a constant factor that may be arbitrarily attached, it is the only invariant of degree $4$.

The equation

\weq{5}{f\,(x_1, x_2, x_3) = 0}

signifies, if $x_1, x_2, x_3$ are coordinates in the plane, a \gesperrt{curve of the fourth order}, which we shall call the \gesperrt{fundamental curve} of the group, or the curve $f$; and there are 168 mappings of the plane onto itself under which to every point of this curve there corresponds a point of the curve. If a point lies on the fundamental curve, then all its image points lie on it as well. The number of distinct image points of a point can sink to a smaller number than 168 — and indeed always to a divisor of 168 — only if the point is a pole.
